\section{重积分 II}

\subsection{重积分的换元法}

\begin{theorem}
	设 $f(x, y)$ 在有界闭区域 $D$ 上可积，若变换 $T: x = x(u, v),\ y = y(u, v)$ 是 $uv$ 平面上分段光滑封闭曲线围成的区域 $\Delta$ 到 $D$ 的双射，且 $x(u, v),\ y(u, v)$ 在 $\Delta$ 有一阶连续偏导数，且
	$$
	J(u, v) = \abs{\frac{\partial(x, y)}{\partial(u, v)}} \neq 0
	$$
	那么
	$$
	\iint_D f(x, y) \dd{x} \dd{y} = \iint_{\Delta} f(x(u, v), y(u, v)) \abs{J(u, v)} \dd{u} \dd{v}
	$$
\end{theorem}

\subsection{极坐标下重积分的计算}

当被积区域是圆形区域（或一部分）或者被积函数是关于 $x^2 + y^2$ 的表达式，那么可以使用极坐标变换化简积分：
$$
T: \begin{dcases}
	x = r \cos \theta \\
	y = r \sin \theta
\end{dcases},\quad 0 \le r < + \infty,\ 0 \le \theta < 2\pi
$$
其 Jacobi 行列式为
$$
J(r, \theta) = \frac{\partial(x, y)}{\partial(r, \theta)} = \begin{vmatrix}
	\cos \theta & -r \sin \theta \\
	\sin \theta & r \cos \theta \\
\end{vmatrix} = r
$$

极坐标变换将 $r \theta$ 平面上的区域 $\Delta$ 映射到 $xy$ 平面上的区域 $D$，那么：
$$
\iint_D f(x, y) \dd{x} \dd{y} = \iint_{\Delta} f(r \cos \theta, r \sin \theta) r \dd{r} \dd{\theta}
$$

\subsection{作业}

\begin{problem}
	习题 16.2.1
	\begin{solution}
		\begin{enumerate}
			\item[\textbf{2)}]
			$$
			\begin{aligned}
				\origin & = \int_{0}^{\pi} \dd{y} \int_{0}^{\pi} \cos (x + y) \dd{x} \\
				& = \int_{0}^{\pi} \ab[\sin (x + y)]_{x = 0}^{\pi} \dd{y} \\
				& = \int_{0}^{\pi} -2 \sin y \dd{y} \\
				& = \ab[2 \cos y]_{y = 0}^{\pi} \\
				& = -4
			\end{aligned}
			$$
		\end{enumerate}
	\end{solution}
\end{problem}

\begin{problem}
	习题 16.2.3
	\begin{solution}
		\begin{enumerate}
			\item[\textbf{1)}]
			$$
			\origin = \int_{0}^{2} \dd{y} \int_{\frac{y}{2}}^{y} f(x, y) \dd{x}
			$$
			\item[\textbf{3)}]
			$$
			\origin = \int_{0}^{a} \dd{y} \int_{\frac{y^2}{2a}}^{a - \sqrt{a^2 - y^2}} f(x, y) \dd{x} + \int_{0}^{a} \dd{y} \int_{a + \sqrt{a^2 - y^2}}^{2a} f(x, y) \dd{x} + \int_{a}^{2a} \dd{y} \int_{\frac{y^2}{2a}}^{2a} f(x, y) \dd{x}
			$$
		\end{enumerate}
	\end{solution}
\end{problem}

\begin{problem}
	习题 16.2.4
	\begin{solution}
		\begin{enumerate}
			\item[\textbf{2)}]
			$$
			\begin{aligned}
				\origin & = \int_{0}^{2} \dd{x} \int_{-\sqrt{2x}}^{\sqrt{2x}} x^2 y^2 \dd{y} \\
				& = \int_{0}^{2} x^2 \cdot \frac{2}{3} (2x)^{\frac{3}{2}} \dd{x} \\
				& = \int_{0}^{2} \frac{4 \sqrt{2}}{3} x^{\frac{7}{2}} \dd{x} \\
				& = \ab[\frac{4 \sqrt{2}}{3} \cdot \frac{2}{9} x^{\frac{9}{2}}]_0^2 = \frac{256}{27}
			\end{aligned}
			$$

			\item[\textbf{4)}] $D$ 是一个菱形区域，因为 $\cos (x + y) = \cos (-x - y)$，所以我们可以只计算 $x > 0$ 的部分：
			$$
			\begin{aligned}
				\origin & = 2 \int_{0}^{1} \dd{x} \int_{x - 1}^{1 - x} \cos (x + y) \dd{y} \\
				& = 2 \int_{0}^{1} \ab[\sin (x + y)]_{x - 1}^{1 - x} \dd{x} \\
				& = 2 \int_{0}^{1} (\sin 1 - \sin (2x - 1)) \dd{x} \\
				& = 2 \ab[x \sin 1 + \frac{1}{2} \cos (2x - 1)]_{0}^{1} = 2 \sin 1
			\end{aligned}
			$$

			\item[\textbf{6)}]
			$$
			\begin{aligned}
				\origin & = \int_{-1}^{0} \dd{x} \int_{x^2}^{1} \ab(y + x f\ab(x^2 + y^2)) \dd{y} + \int_{0}^{1} \dd{x} \int_{x^2}^{1} \ab(y + x f\ab(x^2 + y^2)) \dd{y} \\
				& = \int_{0}^{1} \dd{x} \int_{x^2}^{1} \ab(y - x f\ab(x^2 + y^2)) \dd{x} + \int_{0}^{1} \dd{x} \int_{x^2}^{1} \ab(y + x f\ab(x^2 + y^2)) \dd{y} \\
				& = \int_{0}^{1} \dd{x} \int_{x^2}^{1} 2y \dd{y} \\
				& = \int_{0}^{1} (1 - x^4) \dd{x} = \frac{4}{5}
			\end{aligned}
			$$
		\end{enumerate}
	\end{solution}
\end{problem}

\begin{problem}
	习题 16.2.5
	\begin{proof}
		设 $f(x)$ 的原函数为 $F(x)$ 且 $F(0) = 0$。那么
		$$
		\begin{aligned}
			\lhs & = \int_{0}^{a} f(x) \dd{x} \int_{0}^{x} f(y) \dd{y} \\
			& = \int_{0}^{a} f(x) F(x) \dd{x} \\
			& = \int_{0}^{a} F(x) \dd(F(x)) \\
			& = \frac{1}{2} F(a)^2
		\end{aligned}
		$$
	\end{proof}
\end{problem}

\begin{problem}
	习题 16.2.6
	\begin{solution}
		设极坐标变换将 $r \theta$ 平面下的区域 $\Delta$ 映射到 $D$。
		\begin{enumerate}
			\item[\textbf{1)}] 可知 $\Delta = [0, a] \times [0, 2\pi)$，因此
			$$
			\begin{aligned}
				\origin & = \iint_{\Delta} f(r \cos \theta, r \sin \theta) r \dd{r} \dd{\theta} \\
				& = \int_{0}^{a} \dd{r} \int_{0}^{2\pi} f(r \cos \theta, r \sin \theta) r \dd{\theta}
			\end{aligned}
			$$
			\item[\textbf{2)}] 可知 $0 \le \theta \le \pi, r^2 \le a r \sin \theta \iff r \le a \sin \theta$，因此：
			$$
			\begin{aligned}
				\origin & = \iint_{\Delta} f(r \cos \theta, r \sin \theta) r \dd{r} \dd{\theta} \\
				& = \int_{0}^{\pi} \dd{\theta} \int_{0}^{a \sin \theta} f(r \cos \theta, r \sin \theta) r \dd{r}
			\end{aligned}
			$$
			\item[\textbf{3)}] 可知 $0 \le \theta \le \frac{\pi}{2},\ r \le \frac{a}{\max\{\sin \theta, \cos \theta\}}$，因此：
			$$
			\begin{aligned}
				\origin & = \iint_{\Delta} f(r \cos \theta, r \sin \theta) r \dd{r} \dd{\theta} \\
				& = \int_{0}^{\frac{\pi}{4}} \dd{\theta} \int_{0}^{\frac{a}{\cos \theta}} f(r \cos \theta, r \sin \theta) r \dd{r} + \int_{\frac{\pi}{4}}^{0} \dd{\theta} \int_{0}^{\frac{a}{\sin \theta}} f(r \cos \theta, r \sin \theta) r \dd{r}
			\end{aligned}
			$$
		\end{enumerate}
	\end{solution}
\end{problem}

\begin{problem}
	习题 16.2.7
	\begin{solution}
		\begin{enumerate}
			\item[\textbf{2)}]
			$$
			\begin{aligned}
				\origin & = \int_{1}^{4} \dd{r} \int_{0}^{2\pi} \frac{\sin (\pi r)}{r} \cdot r \dd{\theta} \dd{r} \\
				& = \int_{1}^{4} 2\pi \sin (\pi r) \dd{r} \\
				& = \ab[-2\pi \cos (\pi r)]_{1}^{4} = -2
			\end{aligned}
			$$
			\item[\textbf{3)}]
			$$
			\begin{aligned}
				\origin & = \int_{0}^{\pi} \dd{\theta} \int_{0}^{2 \cos \theta} \sqrt{4 - r^2} r \dd{r} \\
				& = \int_{0}^{\pi} \dd{\theta} \int_{0}^{2 \cos \theta} \frac{1}{2} \sqrt{4 - r^2} \dd(r^2) \\
				& = \int_{0}^{\pi} \ab[-\frac{1}{3} \ab(4 - r^2)^{\frac{3}{2}}]_{0}^{2 \cos \theta} \dd{\theta} \\
				& = \int_{0}^{\pi} \frac{8}{3} \ab(1 - \sin^3 \theta) \dd{\theta} = \frac{8(3 \pi - 4)}{9}
			\end{aligned}
			$$
			\item[\textbf{5)}] 作代换
			$$
			T: \begin{dcases}
				x = a r \cos \theta \\
				y = b r \sin \theta \\
			\end{dcases}\quad,\ 0 \le \theta < 2\pi,\ 0 \le r \le 1
			$$
			Jacobi 行列式为：
			$$
			J(r, \theta) = \begin{vmatrix}
				a \cos \theta & -a r \sin \theta \\
				b \sin \theta & b r \cos \theta \\ 
			\end{vmatrix} = a b r
			$$
			因此
			$$
			\begin{aligned}
				\origin & = \int_{0}^{1} \dd{r} \int_{0}^{2\pi} \sqrt{1 - r^2} r \dd{\theta} \\
				& = 2 \pi \int_{0}^{1} \frac{1}{2} \sqrt{1 - r^2} \dd(r^2) \\
				& = 2 \pi \ab[-\frac{1}{3} \ab(1 - r^2)^{\frac{3}{2}}]_{0}^{1} = \frac{2 \pi}{3}
			\end{aligned}
			$$
		\end{enumerate}
	\end{solution}
\end{problem}

\begin{problem}
	习题 16.2.8
	\begin{solution}
		\begin{enumerate}
			\item[\textbf{1)}] 做变换：
			$$
			T: \begin{dcases}
				x = u - v \\
				y = v \\
			\end{dcases}\quad,\ 0 \le u \le 1,\ 0 \le v \le u
			$$
			那么 Jacobi 行列式
			$$
			J(u, v) = \begin{vmatrix}
				1 & -1 \\
				0 & 1 
			\end{vmatrix} = 1
			$$
			因此
			$$
			\begin{aligned}
				\origin & = \iint_{\Delta} \E^{\frac{v}{u}} \dd{u} \dd{v} \\
				& = \int_{0}^{1} \dd{u} \int_{0}^{u} \E^{\frac{v}{u}} \dd{v} \\
				& = \int_{0}^{1} \E u \dd{u} = \frac{1}{2} \E
			\end{aligned}
			$$

			\item[\textbf{3)}] 做变换：
			$$
			T: \begin{dcases}
				x = u - v \\
				y = v \\
			\end{dcases}\quad,\ 0 \le u \le 1,\ 0 \le v \le u
			$$
			那么 Jacobi 行列式
			$$
			J(u, v) = \begin{vmatrix}
				1 & -1 \\
				0 & 1 
			\end{vmatrix} = 1
			$$
			原限制可以转换为
			$$
			4 \le u \le 12 \land v^2 \le 2 (u - v) \iff 4 \le u \le 12 \land -1 - \sqrt{1 + 2u} \le v \le -1 + \sqrt{1 + 2u}
			$$
			因此
			$$
			\begin{aligned}
				\origin & = \int_{4}^{12} \dd{u} \int_{-1 - \sqrt{1 + 2u}}^{-1 + \sqrt{1 + 2u}} u \dd{v} \\
				& = \int_{4}^{12} 2 u \sqrt{1 + 2u} \dd{u} \\
				& = \int_{3}^{5} (t^2 - 1) t \dd(\frac{t^2 - 1}{2}) \\
				& = \int_{3}^{5} (t^2 - 1) t^2 \dd{t} \\
				& = \ab[\frac{1}{5} t^{5} - \frac{1}{3} t^{3}]_{3}^{5} = \frac{8156}{15}
			\end{aligned}
			$$
		\end{enumerate}
	\end{solution}
\end{problem}

\begin{problem}
	习题 16.2.10
	\begin{solution}
		$$
		\begin{aligned}
			V & = \iint\limits_{x^2 + y^2 \le 1} \sqrt{x^2 + y^2} \dd{x} \dd{y} \\
			& = \iint\limits_{\stackrel{0 \le r \le 1}{0 \le \theta < 2\pi}} r \sqrt{r} \dd{r} \dd{\theta} \\
			& = \int_{0}^{1} \dd{r} \int_{0}^{2\pi} r^{\frac{3}{2}} \dd{\theta} \\
			& = \int_{0}^{1} 2 \pi r^{\frac{3}{2}} \dd{r} = \ab[\frac{4\pi}{5} r^{\frac{5}{2}}]_{0}^{1} = \frac{4\pi}{5}
		\end{aligned}
		$$
	\end{solution}
\end{problem}

\begin{problem}
	习题 16.2.11
	\begin{solution}
		\begin{enumerate}
			\item[\textbf{1)}] 令 $u = x + y, v = \frac{y}{x}$，那么变换
			$$
			T: \begin{dcases}
				x = \frac{u}{1 + v} \\
				y = \frac{uv}{1 + v} \\
			\end{dcases}
			$$
			得到 Jacobi 行列式
			$$
			J(u, v) = \begin{vmatrix}
				\frac{1}{1 + v} & -\frac{u}{(1 + v)^2} \\
				\frac{v}{1 + v} & \frac{u}{(1 + v)^2} \\ 
			\end{vmatrix} = \frac{u}{1 + v}
			$$
			因此，原面积可化为
			$$
			\begin{aligned}
				S & = \iint_D \dd{x} \dd{y} = \int_{a}^{b} \dd{u} \int_{\alpha}^{\beta} \frac{u}{1 + v} \dd{v} \\
				& = \int_{a}^{b} \ab[u \ln (1 + v)]_{v = \alpha}^{\beta} \dd{u} \\
				& = \int_{a}^{b} u \ln \frac{1 + \beta}{1 + \alpha} \dd{u} \\
				& = (b - a) \ln \frac{1 + \beta}{1 + \alpha}
			\end{aligned}
			$$

			\item[\textbf{3)}] 取一个变换
			$$
			T: \begin{dcases}
				u = a_1 x + b_1 y + c_1 \\
				v = a_2 x + b_2 y + c_2 \\
			\end{dcases}
			$$
			那么 Jacobi 行列式为
			$$
			J(u, v) = \begin{vmatrix}
				a_1 & b_1 \\
				a_2 & b_2 
			\end{vmatrix} = a_1 b_2 - a_2 b_1
			$$
			从而得到：
			$$
			\iint\limits_{u^2 + v^2 \le 1} \dd{u} \dd{v} = \iint\limits_{(a_1 x + b_1 y + c_1)^2 + (a_2 x + b_2 y + c_2)^2 \le 1} \abs{a_1 b_2 - a_2 b_1} \dd{x} \dd{y}
			$$
			因此原面积为
			$$
			S = \iint\limits_{(a_1 x + b_1 y + c_1)^2 + (a_2 x + b_2 y + c_2)^2 \le 1} \dd{x} \dd{y} = \frac{\iint_{u^2 + v^2 \le 1} \dd{u} \dd{v}}{\abs{a_1 b_2 - a_2 b_1}} = \frac{\pi}{\abs{a_1 b_2 - a_2 b_1}}
			$$
		\end{enumerate}
	\end{solution}
\end{problem}

\begin{problem}
	习题 16.2.12
	\begin{solution}
		\begin{enumerate}
			\item[\textbf{1)}] 使用极坐标变换：
			$$
			\begin{aligned}
				\origin & = \iint\limits_{\stackrel{0 \le r \le 1}{0 \le \theta < 2\pi}} f(r) r \dd{r} \dd{y} \\
				& = \int_{0}^{1} \dd{r} \int_{0}^{2\pi} f(r) r \dd{\theta} \\
				& = \int_{0}^{1} 2\pi r f(r) \dd{r}
			\end{aligned}
			$$
			\item[\textbf{2)}] 令 $u = \frac{y}{x},\ v = x y$，那么变换为
			$$
			T: \begin{dcases}
				x = \sqrt{\frac{v}{u}} \\
				y = \sqrt{u v}
			\end{dcases}\quad,\ 1 \le  u \le 4,\ 1 \le v \le 2
			$$
			Jacobi 行列式为
			$$
			J(u, v) = \begin{vmatrix}
				2 u^{\frac{1}{2}} v^{\frac{1}{2}} & \frac{2 v^{\frac{3}{2}}}{3 u^{\frac{1}{2}}} \\
				\frac{2}{3} u^{\frac{3}{2}} v^{\frac{1}{2}} & \frac{2}{3} u^{\frac{1}{2}} v^{\frac{3}{2}} 
			\end{vmatrix} = \frac{4}{3} u v^2 - \frac{4}{9} v^2
			$$
			因此
			$$
			\origin = \int_{1}^{4} \dd{u} \int_{1}^{2} f(v) \ab(\frac{4}{3} u v^2 - \frac{4}{9} v^2) \dd{v}
			$$
		\end{enumerate}
	\end{solution}
\end{problem}